\section{Conclusion}
This appendix documented the extension of the \systemname{} framework~\cite{bao2026llmhric} from simulated IAB power allocation to per-S-NSSAI downlink PRB allocation on a running OAI/FlexRIC system, and the implementation and deployment work that extension required. Concretely: an E2SM-RC Style~2 / Action~6 control path and an S-NSSAI-aware downlink allocator were added to OAI; a persistent RC xApp actuates every 100~ms without opening an E2 association per action; 50~ms slice summaries with explicit identity and provenance masks supply the near-RT state; a locally hosted quantized Gemma rApp publishes intent-conditioned targets every 10~s through an A1-like versioned policy service, projected into calibration-derived bounds; and a TD3-style actor--critic learns asynchronously from acknowledged actions and their measured effects, entering the control loop only through validated, atomically swapped Actor snapshots.

\ifFullCampaign
The complete three-scenario, three-arm, five-seed campaign is reported with same-seed paired intervals.
\else
The evidence released with this version is a balanced-load, single-seed pilot of three runs. It shows that the loop closes and is instrumented end to end---all controls acknowledged, five UEs mapped throughout, full effect-window coverage, and every logged transition admissible for learning---and it shows encouraging differences in intent satisfaction and aggregate throughput. It does not establish statistical superiority of any arm, and this paper makes no such claim. The pilot also exposes a specific weakness to address before the full campaign: the guided policy sat on the guidance boundary for the large majority of its actions, so guidance functioned largely as a calibrated constraint rather than as a learning signal.
\fi

What is released is the derived evidence rather than the measurements themselves: the full source and OAI/FlexRIC patches, the campaign specification and validation gates, the result-generation scripts, and, for each executed run, the derived metrics, the generated tables and figures, and a manifest pinning the specification digest and the code revisions. The raw run databases and the trained checkpoints of the pilot were lost and cannot be re-analysed, so every reported number is traceable to the released derived metrics but none can be re-derived from raw data (Section~\ref{sec:limitations}). The remaining work is accordingly to re-run the 45-run campaign with its databases archived on persistent storage, to close the result-generator and validation gaps identified in Section~\ref{sec:limitations}, and to replace every pilot statement with paired multi-scenario estimates. The artifact is available at \codeurlprinted.
